%!TEX root = hems_proj.tex
%%%%%%%%%%%%%%%%%%%%%%%%%%%%%%%%%%%%%%%%%%%%%%%%%%%%%%%%%%%%%%%%%%%%%%%
\chapter{Alkalmazott Algoritmusok}\label{ch:ELM}
%%%%%%%%%%%%%%%%%%%%%%%%%%%%%%%%%%%%%%%%%%%%%%%%%%%%%%%%%%%%%%%%%%%%%%%

\begin{osszefoglal}
Négy metaheurisztikát és egy hibrid eljárást alkalmaztunk: GA, PSO, GWO és GA-PSO hibrid.
\end{osszefoglal}

\section{Genetikus Algoritmus (GA)}
A genetikus algoritmusok alapötletét Holland klasszikus munkája alapozta meg \cite{HollandGA1975}. A módszer fő előnye, hogy jól kezeli a diszkrét változókat és robusztus megoldást nyújt \cite{HollandGA1975}, bár hátránya a lassabb konvergencia. A mérések alapján az algoritmus átlagosan 15--20\%-os költségcsökkenést eredményez.

Az implementált algoritmus működése a következő lépésekből áll. Elsőként generálunk $N$ darab véletlen egyedet, amelyek a kezdeti populációt alkotják, majd minden egyedhez kiszámítjuk a költségfüggvény alapján a fitness értéket. Ezt követi a fő ciklus, amely $G$ generáción keresztül fut. Minden iterációban az elitizmus elve alapján a legjobb két egyed automatikusan továbbjut. A szülőket tournament módszerrel választjuk ki, majd aritmetikai átlagolással és zaj hozzáadásával keresztezzük őket. A sokszínűség fenntartása érdekében Gauss-eloszlású véletlen módosítással mutációt hajtunk végre. Végül az új populáció létrehozása és a határok ellenőrzése után az algoritmus visszatér a legjobb megoldással.

\section{Particle Swarm Optimization (PSO)}
A PSO eredeti formáját Kennedy és Eberhart vezette be \cite{KennedyEberhartPSO1995}.
Sebesség-frissítési egyenlet:
\begin{equation}
V_i^{k+1} = w \cdot V_i^k + c_1 r_1 (P_{best,i} - X_i^k) + c_2 r_2 (G_{best} - X_i^k)
\label{eq:pso_velocity}
\end{equation}
Pozíció-frissítés:
\begin{equation}
X_i^{k+1} = X_i^k + V_i^{k+1}
\label{eq:pso_position}
\end{equation}
Az \eqref{eq:pso_velocity} és \eqref{eq:pso_position} egyenletek adják a PSO alapdinamikáját.
Előny: gyors, egyszerű. Hátrány: lokális optimumra hajlamos. Teljesítmény: $\approx$ 15--18\% költségcsökkenés.

\section{Grey Wolf Optimizer (GWO)}
A GWO algoritmust Mirjalili és szerzőtársai javasolták \cite{MirjaliliGWO2014}.

Hierarchikus struktúra: $\alpha$, $\beta$, $\delta$, $\omega$.
Pozíció-frissítés:
\begin{align}
\vec{D}_\alpha &= |\vec{C}_1 \cdot \vec{X}_\alpha - \vec{X}| \\
\vec{X}_1 &= \vec{X}_\alpha - \vec{A}_1 \cdot \vec{D}_\alpha \\
\vec{X}(t+1) &= \frac{\vec{X}_1 + \vec{X}_2 + \vec{X}_3}{3}
\end{align}
Előny: kevesebb paraméter, jobb exploration-exploitation egyensúly.
Teljesítmény: $\approx$ 18--20\% költségcsökkenés.

\section{Hibrid GA-PSO}
A GA és PSO kombinálása ismert hibridizációs irány \cite{KumarGA}.

A hibridizáció logikája a két módszer erősségeit ötvözi: a keresés első fázisában a PSO végzi a gyors feltérképezést a sebesség-frissítési egyenletek segítségével. Ezt követően a GA veszi át az irányítást, amely keresztezéssel és mutációval biztosítja a megoldások diverzifikációját. A két fázis közötti váltás adaptív módon történik, így az algoritmus képes elkerülni a korai konvergenciát, elérve ezzel mintegy 18--22\%-os költségcsökkenést.